% =============================================================================
% Presentación Beamer P04 – LDA / Perfil de Velocidades
% Técnicas de Medida – Ingeniería Aeronáutica
% Compilar con: latexmk -xelatex -interaction=nonstopmode presentacion_P04.tex
% =============================================================================
\documentclass[aspectratio=169,12pt]{beamer}
\usetheme{Madrid}
\usecolortheme{seahorse}
\usefonttheme{professionalfonts}

% ---- Encoding & Language ----
\usepackage{fontspec}
\usepackage{csquotes}
\usepackage[spanish]{babel}

% ---- Math & Units ----
\usepackage{amsmath,amssymb}
\usepackage{siunitx}
\sisetup{output-decimal-marker={,}}

% ---- TikZ (babel first!) ----
\usepackage{tikz}
\usetikzlibrary{babel}
\usetikzlibrary{arrows.meta,positioning,calc,shapes.geometric,
                decorations.markings,decorations.pathmorphing,patterns}

% ---- Tables & Graphics ----
\usepackage{booktabs}
\usepackage{colortbl}
\usepackage{tabularx}
\usepackage{array}
\usepackage{graphicx}

% ---- Code Listings ----
\usepackage{listings}
\lstset{
  basicstyle=\ttfamily\tiny,
  keywordstyle=\color{blue!80}\bfseries,
  commentstyle=\color{green!50!black},
  backgroundcolor=\color{gray!10},
  frame=single,
  breaklines=true
}

% ---- Custom Colors ----
\definecolor{aeroblue}{RGB}{0,51,102}
\definecolor{aerored}{RGB}{153,0,0}
\definecolor{aerogreen}{RGB}{0,102,51}
\definecolor{aeroyellow}{RGB}{204,153,0}
\definecolor{aerolaser}{RGB}{0,180,0}
\definecolor{aerogreendark}{RGB}{0,80,40}
\definecolor{aerogreenlight}{RGB}{0,140,70}

% ---- Beamer Color Theme (GREEN) ----
\setbeamercolor{frametitle}{bg=aerogreen,fg=white}
\setbeamercolor{title}{bg=aerogreen,fg=white}
\setbeamercolor{structure}{fg=aerogreen}
\setbeamercolor{palette primary}{bg=aerogreen,fg=white}
\setbeamercolor{palette secondary}{bg=aerogreenlight,fg=white}
\setbeamercolor{palette tertiary}{bg=aerogreendark,fg=white}
\setbeamercolor{block title}{bg=aerogreen,fg=white}
\setbeamercolor{block body}{bg=aerogreen!8}
\setbeamercolor{block title alerted}{bg=aerored,fg=white}
\setbeamercolor{block body alerted}{bg=aerored!8}
\setbeamercolor{block title example}{bg=aeroblue,fg=white}
\setbeamercolor{block body example}{bg=aeroblue!8}
\setbeamercolor{item}{fg=aerogreen}
\setbeamercolor{subitem}{fg=aerogreenlight}

% ---- Image path ----
\graphicspath{{../../../Practicas/P04_LDA_Perfil_Velocidad/data/images/}}

% ---- Title ----
\title[P04 -- LDA]{P04: Anemometría Láser Doppler\\Perfil de Velocidades}
\subtitle{Técnicas de Medida -- Ingeniería Aeronáutica}
\author{Laboratorio de Técnicas de Medida}
\institute{Departamento de Ingeniería Aeroespacial}
\date{\today}

% =============================================================================
\begin{document}
% =============================================================================

\begin{frame}
	\titlepage
\end{frame}

\begin{frame}{Contenido}
	\tableofcontents
\end{frame}

% =============================================================================
\section{Introducción}
% =============================================================================

\begin{frame}{¿Qué es la Anemometría Láser Doppler (LDA)?}
	\begin{columns}
		\begin{column}{0.5\textwidth}
			\begin{itemize}\small
				\item Técnica \textbf{óptica no intrusiva} para
				      medir velocidad puntual en fluidos
				\item Basada en el \textbf{efecto Doppler} de luz
				      dispersada por partículas trazadoras
				\item Dos haces láser coherentes $\to$ patrón
				      de franjas $\to$ $f_D \propto U$
				\item También denominada \textbf{LDV}
				      (\emph{Laser Doppler Velocimetry})
			\end{itemize}
			\vspace{0.3em}
			\begin{alertblock}{Ventaja clave}
				\textbf{Auto-calibrante}: $U = f_D \cdot \delta_f$\\
				depende solo de la geometría óptica
			\end{alertblock}
		\end{column}
		\begin{column}{0.47\textwidth}
			\begin{tikzpicture}[scale=0.75,>=Stealth]
				% Two laser beams
				\draw[aerolaser,ultra thick] (0,1.5) -- (3.0,0);
				\draw[aerolaser,ultra thick] (0,-1.5) -- (3.0,0);
				% Measurement volume
				\fill[aerolaser!25] (3.0,0) ellipse (0.25 and 0.12);
				\draw[aerolaser,thick] (3.0,0) ellipse (0.25 and 0.12);
				% Fringes
				\foreach \dx in {-0.18,-0.09,0.0,0.09,0.18} {
						\draw[aerolaser!50,very thin]
						(3+\dx,-0.12) -- (3+\dx,0.12);
					}
				% Particle
				\fill[gray!80] (3.0,0) circle (0.06);
				% Flow arrow
				\draw[->,blue!60,thick] (2.2,-0.3) -- (3.8,-0.3)
				node[right,font=\tiny] {$U$};
				% Scattered light
				\draw[gray,dashed,->] (3.0,0.12) .. controls (3.5,1.0)
				.. (4.5,1.0) node[right,font=\tiny] {$f_D$};
				% Beam labels
				\node[aerolaser,above,font=\tiny] at (1.2,0.9)
				{Haz 1};
				\node[aerolaser,below,font=\tiny] at (1.2,-0.9)
				{Haz 2};
				% Fringe spacing
				\draw[<->,aerolaser,thin] (2.78,-0.30) -- (3.22,-0.30)
				node[midway,below,font=\tiny] {$\delta_f$};
				% Laser box
				\draw[fill=aerogreen!20,draw=aerogreen]
				(-1,-0.15) rectangle (0,0.15);
				\node[font=\tiny] at (-0.5,0) {Láser};
				\draw[->,aerolaser] (0,0.07) -- (0.3,1.1);
				\draw[->,aerolaser] (0,-0.07) -- (0.3,-1.1);
				% PMT
				\draw[fill=gray!30,draw=black] (4.5,0.7) rectangle (5.3,1.3);
				\node[font=\tiny] at (4.9,1.0) {PMT};
				% BSA
				\draw[fill=aerogreen!20,draw=aerogreen]
				(5.5,0.7) rectangle (6.4,1.3);
				\node[font=\tiny,aerogreen] at (5.95,1.0) {BSA};
				\draw[->] (5.3,1.0) -- (5.5,1.0);
			\end{tikzpicture}
		\end{column}
	\end{columns}
\end{frame}

\begin{frame}{Aplicaciones y Ventajas vs.\ HWA}
	\begin{columns}
		\begin{column}{0.5\textwidth}
			\begin{enumerate}\small
				\item Perfiles de velocidad en túneles de viento
				\item Estelas de perfiles aerodinámicos ($C_D$)
				\item Aerodinámica interna de motores
				\item Flujo en rotores de helicóptero
				\item Validación de CFD
			\end{enumerate}
			\vspace{0.3em}
			\begin{exampleblock}{Ventajas vs.\ HWA (P03)}
				\begin{itemize}\small
					\item \textbf{No intrusiva} (sin sonda)
					\item Mide \textbf{flujos inversos} (Bragg)
					\item \textbf{No requiere calibración}
					\item Opera en alta temperatura
				\end{itemize}
			\end{exampleblock}
		\end{column}
		\begin{column}{0.47\textwidth}
			\begin{alertblock}{Limitaciones}
				\begin{itemize}\small
					\item Necesita acceso óptico (ventanas)
					\item Requiere partículas trazadoras
					      (\emph{seeding})
					\item El fluido debe ser transparente
					\item Coste y complejidad elevados
					\item Medida puntual (no campo completo)
				\end{itemize}
			\end{alertblock}
			\vspace{0.3em}
			\centering\scriptsize
			\begin{tabular}{lcc}
				\toprule
				           & \textbf{HWA} & \textbf{LDA} \\
				\midrule
				Intrusiva  & Sí           & No           \\
				Calibr.    & Sí           & No           \\
				Res.\ temp & Muy alta     & Alta         \\
				Coste      & Medio        & Alto         \\
				\bottomrule
			\end{tabular}
		\end{column}
	\end{columns}
\end{frame}

\begin{frame}{Objetivos de Aprendizaje}
	\begin{enumerate}
		\item Comprender el efecto Doppler aplicado a la dispersión de
		      luz por partículas trazadoras
		\item Describir el modelo de franjas y calcular $\delta_f$
		      del volumen de medida
		\item Relacionar $f_D$ con $U$: $U = f_D \cdot \delta_f$
		\item Calcular $d_{mv}$, $l_{mv}$ y $N_f$ a partir
		      de los parámetros ópticos
		\item Explicar la función de la celda de Bragg ($f_s$)
		\item Procesar datos LDA del laboratorio (\texttt{files/P4/})
		      para obtener $U(y)$
		\item Aplicar filtrado IQR y calcular $\bar{U}$, $\sigma_U$,
		      $I_u$ por posición
		\item Propagar incertidumbre $f_D \to U$ por método GUM
		\item Realizar evaluación de riesgos con la matriz $5\times5$
	\end{enumerate}
\end{frame}

% =============================================================================
\section{Fundamentos Teóricos}
% =============================================================================

\begin{frame}{Efecto Doppler y Modelo de Franjas}
	\begin{columns}
		\begin{column}{0.5\textwidth}
			\textbf{Frecuencia Doppler:}
			\begin{equation*}
				f_D = \frac{2\sin(\theta/2)}{\lambda}\,U
			\end{equation*}

			\textbf{Separación de franjas:}
			\begin{equation*}
				\boxed{\delta_f = \frac{\lambda}{2\sin(\theta/2)}}
			\end{equation*}

			\textbf{Velocidad de la partícula:}
			\begin{equation*}
				\boxed{U = f_D \cdot \delta_f}
			\end{equation*}

			\vspace{0.3em}
			La relación es \textbf{lineal}: no se necesita
			calibración, solo conocer $\lambda$ y $\theta$.
		\end{column}
		\begin{column}{0.47\textwidth}
			\begin{tikzpicture}[scale=0.8,>=Stealth]
				% Two beams
				\draw[aerolaser,ultra thick] (-2.5,1.5) -- (0,0);
				\draw[aerolaser,ultra thick] (-2.5,-1.5) -- (0,0);
				% Volume ellipse
				\fill[aerolaser!15] (0,0) ellipse (0.5 and 0.22);
				\draw[aerolaser,thick,dashed] (0,0) ellipse (0.5 and 0.22);
				% Fringes
				\foreach \dx in {-0.4,-0.3,-0.2,-0.1,0.0,0.1,0.2,0.3,0.4} {
						\draw[aerolaser!40,thin] (\dx,-0.22) -- (\dx,0.22);
					}
				% Fringe spacing
				\draw[<->,aerored,thick] (-0.1,-0.32) -- (0.1,-0.32)
				node[midway,below,font=\tiny] {$\delta_f$};
				% d_mv
				\draw[<->,aerogreen,thick] (-0.5,-0.5) -- (0.5,-0.5)
				node[midway,below,font=\tiny] {$d_{mv}$};
				% Particle
				\fill[blue!80] (0,0) circle (0.06);
				\draw[->,blue!60,thick] (-0.7,0) -- (0.7,0)
				node[right,font=\tiny] {$U_\perp$};
				% Angle
				\draw[gray,thin] (-2.5,0) -- (0,0);
				\draw[gray,thin] (-1.5,0) arc (0:31:1.0);
				\node[gray,font=\tiny] at (-1.2,0.22) {$\theta/2$};
				% Scattered light
				\draw[gray!60,dashed,thick,->] (0,0.22)
				.. controls (0.5,1.2) .. (1.8,1.2)
				node[right,font=\tiny] {Luz dispersada};
				% Labels
				\node[aerolaser,above,font=\tiny] at (-1.5,1.0) {Haz 1};
				\node[aerolaser,below,font=\tiny] at (-1.5,-1.0) {Haz 2};
			\end{tikzpicture}
		\end{column}
	\end{columns}
\end{frame}

\begin{frame}{Volumen de Medida y Número de Franjas}
	\begin{columns}
		\begin{column}{0.5\textwidth}
			\textbf{Haces gaussianos enfocados:}
			\begin{align*}
				d_{mv} & = \frac{4\,f_L\,\lambda}{\pi\,d_b}
				\quad\text{(diámetro)}                      \\
				l_{mv} & = \frac{d_{mv}}{\sin(\theta/2)}
				\quad\text{(longitud)}                      \\
				N_f    & = \frac{d_{mv}}{\delta_f}
				\quad\text{(Nº franjas)}
			\end{align*}

			\vspace{0.3em}
			\begin{block}{Parámetros del laboratorio}
				\begin{itemize}\scriptsize
					\item $\lambda = \SI{514,5}{\nm}$ (Ar-Ion)
					\item $\theta/2 = \SI{4,5}{\degree}$
					\item $f_L = \SI{250}{\mm}$
					\item $d_b = \SI{1,2}{\mm}$
					\item $f_s = \SI{40}{\MHz}$ (Bragg)
				\end{itemize}
			\end{block}
		\end{column}
		\begin{column}{0.47\textwidth}
			\textbf{Cálculos numéricos:}
			\begin{align*}
				\delta_f & = \frac{514{,}5\!\times\!10^{-9}}{2\sin(4{,}5°)}      \\
				         & = \SI{3,279}{\micro\metre}                            \\[0.5em]
				d_{mv}   & = \frac{4 \times 250 \times 514{,}5\!\times\!10^{-9}}
				{\pi\times 1{,}2\!\times\!10^{-3}}                               \\
				         & \approx \SI{136,5}{\micro\metre}                      \\[0.5em]
				l_{mv}   & = \frac{136{,}5}{\sin(4{,}5°)}
				\approx \SI{1,74}{\mm}                                           \\[0.5em]
				N_f      & = 136{,}5 / 3{,}279 \approx 42
			\end{align*}

			\begin{alertblock}{}
				\centering
				$u_U \approx \SI{0,019}{\m\per\s}$\\
				($\approx 0{,}95\%$ a \SI{2}{\m\per\s})
			\end{alertblock}
		\end{column}
	\end{columns}
\end{frame}

\begin{frame}{Celda de Bragg -- Desplazamiento de Frecuencia}
	\begin{columns}
		\begin{column}{0.5\textwidth}
			\textbf{Problema:} sin Bragg, $f_D > 0$ tanto si
			$U > 0$ como si $U < 0$ (ambigüedad).

			\vspace{0.5em}
			\textbf{Solución:} desplazar uno de los haces
			$f_s = \SI{40}{\MHz}$:
			\begin{equation*}
				\boxed{f_\mathrm{obs} = f_s + f_D}
			\end{equation*}

			\begin{itemize}\small
				\item $U > 0 \Rightarrow f_\mathrm{obs} > f_s$
				\item $U < 0 \Rightarrow f_\mathrm{obs} < f_s$
				\item Permite medir \textbf{recirculación}
			\end{itemize}
		\end{column}
		\begin{column}{0.47\textwidth}
			\begin{tikzpicture}[scale=0.65,>=Stealth]
				% Frequency axis
				\draw[->] (0,0) -- (6.5,0)
				node[right,font=\tiny] {$f$};
				\draw[->] (0,0) -- (0,4)
				node[above,font=\tiny] {Amplitud};
				% f_s
				\draw[aerogreen,very thick] (3,0) -- (3,3.5);
				\node[aerogreen,below,font=\tiny] at (3,0) {$f_s$};
				% f_D positive
				\draw[aerogreen,very thick] (4.5,0) -- (4.5,2.5);
				\node[aerogreen,below,font=\tiny] at (4.5,0)
				{$f_s\!+\!f_D$};
				\draw[<->,aerogreen,thick] (3,3.7) -- (4.5,3.7)
				node[midway,above,font=\tiny] {$f_D > 0$};
				% f_D negative
				\draw[aerored,very thick] (1.5,0) -- (1.5,2.0);
				\node[aerored,below,font=\tiny] at (1.5,0)
				{$f_s\!-\!|f_D|$};
				\draw[<->,aerored,thick] (1.5,2.5) -- (3,2.5)
				node[midway,above,font=\tiny] {$f_D < 0$};
			\end{tikzpicture}

			\vspace{0.3em}
			\begin{exampleblock}{Modulador acusto-óptico}
				\scriptsize
				Cristal piezoeléctrico genera onda acústica
				que difracta un haz, desplazándolo $f_s$.
			\end{exampleblock}
		\end{column}
	\end{columns}
\end{frame}

\begin{frame}{Señal de Ráfaga (\emph{Burst}) y Seeding}
	\begin{columns}
		\begin{column}{0.5\textwidth}
			\textbf{Señal \emph{burst}:}\\
			Cada partícula genera una ráfaga sinusoidal
			modulada por la envolvente gaussiana del haz.

			\vspace{0.3em}
			\begin{tikzpicture}[xscale=1.2,yscale=0.7]
				\draw[->] (0,-1.2) -- (5,0) node[right,font=\tiny] {$t$};
				\draw[->] (0,-1.2) -- (0,1.6)
				node[above,font=\tiny] {PMT};
				% Envelope
				\draw[aerored,thick,dashed]
				plot[domain=0.5:4.5,samples=50]
				(\x,{1.1*exp(-0.4*(\x-2.5)^2)});
				\draw[aerored,thick,dashed]
				plot[domain=0.5:4.5,samples=50]
				(\x,{-1.1*exp(-0.4*(\x-2.5)^2)});
				% Burst signal
				\draw[aerogreen,thick]
				plot[domain=0.5:4.5,samples=150]
				(\x,{1.1*exp(-0.4*(\x-2.5)^2)*sin(deg(10*\x))});
				% Period
				\draw[<->,aerogreen,thick] (2.15,-0.9) -- (2.47,-0.9)
				node[midway,below,font=\tiny] {$1/f_D$};
			\end{tikzpicture}
		\end{column}
		\begin{column}{0.47\textwidth}
			\textbf{Partículas trazadoras:}
			\begin{itemize}\small
				\item $d_p \approx \SI{1}{\micro\m}$
				      (aceite nebulizado)
				\item $\rho_p \approx \SI{900}{\kg\per\m\cubed}$
				\item Tiempo de respuesta:
				      \begin{equation*}
					      \tau_p = \frac{\rho_p\,d_p^2}{18\,\mu}
					      \approx \SI{2,8}{\micro\s}
				      \end{equation*}
				\item Sigue turbulencia hasta
				      ${\sim}\SI{50}{\kHz}$
			\end{itemize}

			\vspace{0.3em}
			\begin{alertblock}{Criterio}
				$>\!1000$ ráfagas/s por punto\\
				BSA extrae $f_D$ por FFT
			\end{alertblock}
		\end{column}
	\end{columns}
\end{frame}

\begin{frame}{Configuración Óptica del Sistema LDA}
	\centering
	\begin{tikzpicture}[>=Stealth,scale=0.78]
		% Laser source
		\draw[fill=aerogreen!20,draw=aerogreen,very thick]
		(-1,-0.35) rectangle (1,0.35);
		\node[font=\tiny\bfseries,align=center] at (0,0)
		{Ar-Ion\\\SI{514,5}{\nm}};
		% Beam splitter
		\draw[fill=gray!30,draw=black,thick]
		(1.8,-0.35) rectangle (2.0,0.35);
		\node[below,font=\tiny] at (1.9,-0.45) {Divisor};
		\draw[aerolaser,very thick,->] (1,0) -- (1.8,0);
		% Bragg cell
		\draw[fill=yellow!20,draw=black,thick]
		(2.6,-0.35) rectangle (3.0,0.35);
		\node[below,font=\tiny] at (2.8,-0.45) {Bragg};
		\draw[aerolaser,very thick] (2.0,0.08) -- (2.6,0.08);
		\draw[aerolaser,very thick] (2.0,-0.08) -- (2.6,-0.08);
		% Two beams diverging
		\draw[aerolaser,very thick] (3.0,0.08) -- (4.3,0.8);
		\draw[aerolaser,very thick] (3.0,-0.08) -- (4.3,-0.8);
		% Focusing lens
		\draw[fill=cyan!20,draw=black,thick]
		(4.3,-1.0) -- (4.5,0) -- (4.3,1.0) -- (4.1,1.0)
		-- (4.3,0) -- (4.1,-1.0) -- cycle;
		\node[below,font=\tiny] at (4.3,-1.1) {$f_L$};
		% Measurement volume
		\fill[aerogreen!30] (5.8,0) ellipse (0.12 and 0.06);
		\draw[aerogreen,thick] (5.8,0) ellipse (0.12 and 0.06);
		% Fringes
		\foreach \dx in {-0.08,-0.04,0.0,0.04,0.08} {
				\draw[aerolaser!50,very thin]
				(5.8+\dx,-0.06) -- (5.8+\dx,0.06);
			}
		% Particle
		\draw[->,blue!60,thick] (5.2,-0.3) -- (6.4,-0.3)
		node[right,font=\tiny] {Partícula};
		% Scattered light
		\draw[gray,dashed,->] (5.8,0.06)
		.. controls (6.3,0.8) .. (7.5,0.8);
		% Receiving optics
		\draw[fill=cyan!20,draw=black,thick]
		(7.5,0.5) -- (7.7,0.8) -- (7.5,1.1) -- (7.3,1.1)
		-- (7.5,0.8) -- (7.3,0.5) -- cycle;
		\node[above,font=\tiny] at (7.5,1.2) {Lente rec.};
		% PMT
		\draw[fill=gray!30,draw=black]
		(8.0,0.5) rectangle (8.8,1.1);
		\node[font=\tiny] at (8.4,0.8) {PMT};
		\draw[->,gray!50] (7.5,0.8) -- (8.0,0.8);
		% BSA
		\draw[fill=aerogreen!20,draw=aerogreen]
		(9.1,0.5) rectangle (10.1,1.1);
		\node[font=\tiny,aerogreen] at (9.6,0.8) {BSA};
		\draw[->] (8.8,0.8) -- (9.1,0.8);
		% Output
		\draw[->,aerogreen,thick] (10.1,0.8) -- (11.0,0.8)
		node[right,font=\tiny] {$U = f_D\,\delta_f$};
		% Labels
		\node[above,font=\tiny,aerolaser] at (3.8,0.6) {Haz 1};
		\node[below,font=\tiny,aerolaser] at (3.8,-0.6) {Haz 2};
		\node[below,font=\tiny,aerolaser] at (5.8,-0.12) {Vol.\ med.};
		% Angle
		\draw[gray,thin] (4.3,0) -- (5.8,0);
		\draw[gray,thin] (5.3,0) arc (0:28:0.5);
		\node[gray,font=\tiny] at (5.4,0.13) {$\theta/2$};
	\end{tikzpicture}

	\vspace{0.5em}
	\begin{tikzpicture}
		\node[draw=aerogreen,thick,rounded corners,fill=aerogreen!8,
			text width=12cm,align=center,font=\scriptsize] {
			\textbf{Cadena:} Láser Ar-Ion $\to$ Divisor $\to$
			Bragg ($f_s = \SI{40}{\MHz}$) $\to$ Lente ($f_L = \SI{250}{\mm}$)
			$\to$ Volumen de medida ($\delta_f = \SI{3,28}{\micro\m}$)
			$\to$ PMT $\to$ BSA $\to$ $f_D$ $\to$ $U$
		};
	\end{tikzpicture}
\end{frame}

% =============================================================================
\section{Datos Experimentales}
% =============================================================================

\begin{frame}{Estructura de Datos -- \texttt{files/P4/}}
	\begin{columns}
		\begin{column}{0.45\textwidth}
			\centering\scriptsize
			\begin{tabular}{lccr}
				\toprule
				\textbf{Carpeta} & \textbf{Arch.}
				                 & \textbf{Lín./arch.}
				                 & \textbf{Total}                                   \\
				\midrule
				FX01G00          & 15                  & 2\,000 & 30\,000           \\
				FX02G00          & 15                  & 1\,138 & 17\,070           \\
				FX03G00          & 15                  & 2\,000 & 30\,000           \\
				FX04G00          & 15                  & 2\,000 & 30\,000           \\
				\midrule
				\textbf{Total}   &                     &        & \textbf{107\,070} \\
				\bottomrule
			\end{tabular}

			\vspace{0.3em}
			\textit{\tiny FX02G00: 1\,138 muestras/archivo\\
				(anomalía de adquisición)}
		\end{column}
		\begin{column}{0.52\textwidth}
			\begin{block}{Formato de datos (5 columnas)}
				\scriptsize
				\begin{tabular}{cl}
					\toprule
					\textbf{Col.} & \textbf{Contenido}           \\
					\midrule
					1             & Índice de muestra            \\
					2             & Tiempo de llegada (ms)       \\
					3             & SNR / tiem.\ tránsito        \\
					4             & \textbf{Velocidad $U$} (m/s) \\
					5             & \textbf{Velocidad $V$} (m/s) \\
					\bottomrule
				\end{tabular}
			\end{block}
			\vspace{0.2em}
			\begin{exampleblock}{Lectura automática}
				\scriptsize
				\texttt{pd.read\_csv(file, sep=r'\textbackslash s+')}\\
				$U$ = Col.~4 (principal), $V$ = Col.~5 (transversal)
			\end{exampleblock}
		\end{column}
	\end{columns}
\end{frame}

\begin{frame}{Estadísticos Verificados por Carpeta}
	\centering\small
	\begin{tabular}{l S[table-format=5.0]
			S[table-format=1.4] S[table-format=1.4]
			S[table-format=2.2] S[table-format=1.5]}
		\toprule
		\textbf{Carpeta}
		        & {$N$}              & {$\bar{U}$ (m/s)}
		        & {$\sigma_U$ (m/s)}
		        & {$I_u$ (\%)}
		        & {$\bar{V}$ (m/s)}                                                 \\
		\midrule
		FX01G00 & 30000              & 0.5086            & 0.0731 & 14.37 & 0.01422 \\
		FX02G00 & 29138              & 1.0133            & 0.1146 & 11.31 & 0.02078 \\
		FX03G00 & 30000              & 1.9633            & 0.2191 & 11.16 & 0.07125 \\
		FX04G00 & 30000              & 2.9255            & 0.3486 & 11.92 & 0.07639 \\
		\bottomrule
	\end{tabular}

	\vspace{0.5em}
	\begin{tikzpicture}
		\node[draw=aerogreen,thick,rounded corners,fill=aerogreen!8,
			text width=11cm,align=center,font=\scriptsize] {
			$\bar{U}$ crece de $\sim$\SI{0,5}{\m\per\s} a
			$\sim$\SI{2,9}{\m\per\s} $|$
			$I_u \approx 11$--$14\%$ (flujo moderadamente turbulento) $|$
			$\bar{V}/\bar{U} < 3\%$ (flujo unidimensional)
		};
	\end{tikzpicture}
\end{frame}

\begin{frame}{Parámetros del Sistema LDA}
	\centering
	\begin{tabular}{llr}
		\toprule
		\textbf{Parámetro}            & \textbf{Valor}
		                              & \textbf{Unidad}        \\
		\midrule
		Tipo láser                    & Ar-Ion            & -- \\
		Longitud de onda $\lambda$    & 514,5
		                              & \si{\nm}               \\
		Semi-ángulo $\theta/2$        & 4,5
		                              & \si{\degree}           \\
		Focal lente transmisora $f_L$ & 250
		                              & \si{\mm}               \\
		Diámetro del haz $d_b$        & 1,2
		                              & \si{\mm}               \\
		Desplazamiento Bragg $f_s$    & 40
		                              & \si{\MHz}              \\
		Separación de franjas         & 3,279
		                              & \si{\micro\metre}      \\
		Diámetro vol.\ medida         & 136,5
		                              & \si{\micro\metre}      \\
		Longitud vol.\ medida         & 1,74
		                              & \si{\mm}               \\
		Nº de franjas $N_f$           & 42
		                              & --                     \\
		Carpetas de datos             & 4 (FX01--FX04)
		                              & --                     \\
		Muestras por archivo          & 2\,000 (típico)
		                              & --                     \\
		\bottomrule
	\end{tabular}
\end{frame}

% =============================================================================
\section{Procesamiento de Datos}
% =============================================================================

\begin{frame}{Cadena de Procesamiento (5 Pasos)}
	\centering
	\begin{tikzpicture}[node distance=1.0cm,>=Stealth,
			block/.style={rectangle,draw=aerogreen,thick,
					fill=aerogreen!8,text width=4.5cm,align=center,
					minimum height=0.75cm,font=\scriptsize,
					rounded corners=3pt}]
		\node[block] (read) {\textbf{Paso 1:}
			Leer todos los \texttt{.txt}\\
			de la carpeta (15 archivos)};
		\node[block,below=of read] (concat) {\textbf{Paso 2:}
			Concatenar $U$ y $V$\\
			de las 15 iteraciones};
		\node[block,below=of concat] (iqr) {\textbf{Paso 3:}
			Filtrado IQR ($k = 1{,}5$)\\
			Eliminar outliers};
		\node[block,below=of iqr] (stats) {\textbf{Paso 4:}
			Calcular $\bar{U}$, $\sigma_U$,\\
			$I_u$, $\bar{V}$, SE};
		\node[block,below=of stats,draw=aerored,fill=aerored!8]
		(profile) {\textbf{Paso 5:}
			Construir perfil $U(y)$\\
			con barras de error};
		% Arrows
		\foreach \a/\b in {read/concat,concat/iqr,iqr/stats,
				stats/profile} {
				\draw[->,aerogreen,thick] (\a) -- (\b);
			}
		% Loop
		\draw[aerored,thick,->] (stats.east) -- ++(1.2,0)
		|- (read.east)
		node[pos=0.25,right,font=\tiny,aerored] {$\times 4$ carpetas};
		% Counts
		\node[right=0.3cm,font=\tiny,gray] at (read.west)
		{$\sim$2\,000 lín./arch.};
		\node[right=0.3cm,font=\tiny,gray] at (concat.west)
		{$N \sim 30\,000$};
		\node[right=0.3cm,font=\tiny,gray] at (iqr.west)
		{3--6\% eliminados};
	\end{tikzpicture}
\end{frame}

\begin{frame}{Filtrado IQR de Outliers}
	\begin{columns}
		\begin{column}{0.5\textwidth}
			\textbf{Rango intercuartílico:}
			\begin{equation*}
				\mathrm{IQR} = Q_3 - Q_1
			\end{equation*}

			\textbf{Límites de aceptación:}
			\begin{equation*}
				\bigl[Q_1 - 1{,}5\,\mathrm{IQR},\;
					Q_3 + 1{,}5\,\mathrm{IQR}\bigr]
			\end{equation*}

			\vspace{0.3em}
			\begin{exampleblock}{Resultados del filtrado}
				\scriptsize
				\begin{tabular}{lr}
					\toprule
					\textbf{Carpeta} & \textbf{Outliers (\%)} \\
					\midrule
					FX01G00          & 3,2\%                  \\
					FX02G00          & 3,3\%                  \\
					FX03G00          & 5,7\%                  \\
					FX04G00          & 5,7\%                  \\
					\bottomrule
				\end{tabular}
			\end{exampleblock}
		\end{column}
		\begin{column}{0.47\textwidth}
			\begin{tikzpicture}[scale=0.6,>=Stealth]
				% Box plot conceptual
				\draw[aerogreen,thick] (0,0) rectangle (3.5,2);
				\draw[aerogreen,very thick] (0,1.2) -- (3.5,1.2);
				\draw[aerogreen,thick] (1.75,0) -- (1.75,-0.8);
				\draw[aerogreen,thick] (1.75,2) -- (1.75,3.0);
				\draw[aerogreen,thick] (0.8,-0.8) -- (2.7,-0.8);
				\draw[aerogreen,thick] (0.8,3.0) -- (2.7,3.0);
				% Outlier threshold
				\draw[aerored,thick,dashed]
				(-0.3,-2.0) -- (3.8,-2.0);
				\draw[aerored,thick,dashed]
				(-0.3,4.2) -- (3.8,4.2);
				% Labels
				\node[left,font=\tiny] at (0,0) {$Q_1$};
				\node[left,font=\tiny] at (0,1.2) {Med};
				\node[left,font=\tiny] at (0,2) {$Q_3$};
				\node[left,font=\tiny,aerored] at (-0.3,-2.0)
				{$Q_1\!-\!1{,}5\,\mathrm{IQR}$};
				\node[left,font=\tiny,aerored] at (-0.3,4.2)
				{$Q_3\!+\!1{,}5\,\mathrm{IQR}$};
				% Outlier points
				\fill[aerored] (1.75,-2.5) circle (0.1);
				\fill[aerored] (1.2,4.8) circle (0.1);
				\fill[aerored] (2.5,5.0) circle (0.1);
				\node[aerored,right,font=\tiny] at (2.7,4.9) {Outliers};
				% IQR brace
				\draw[<->,thick] (3.8,0) -- (3.8,2)
				node[midway,right,font=\tiny] {IQR};
			\end{tikzpicture}

			\vspace{0.3em}
			\begin{alertblock}{Ventaja IQR}
				Más \textbf{robusto} que $\pm2\sigma$:\\
				insensible a los propios outliers
			\end{alertblock}
		\end{column}
	\end{columns}
\end{frame}

\begin{frame}{Estadísticos por Posición}
	\begin{columns}
		\begin{column}{0.5\textwidth}
			\textbf{Velocidad media:}
			\begin{equation*}
				\bar{U}(y_k) = \frac{1}{N_k}\sum_{i=1}^{N_k} U_i
			\end{equation*}

			\textbf{Desviación estándar:}
			\begin{equation*}
				\sigma_U(y_k) = \sqrt{\frac{1}{N_k\!-\!1}
					\sum_{i=1}^{N_k}(U_i - \bar{U})^2}
			\end{equation*}

			\textbf{Intensidad de turbulencia:}
			\begin{equation*}
				\boxed{I_u = \frac{\sigma_U}{\bar{U}} \times 100\%}
			\end{equation*}

			\textbf{Error estándar:}
			\begin{equation*}
				\mathrm{SE} = \frac{\sigma_U}{\sqrt{N_k}}
			\end{equation*}
		\end{column}
		\begin{column}{0.47\textwidth}
			\begin{tikzpicture}[scale=0.6]
				\draw[->] (0,0) -- (5.5,0)
				node[right,font=\tiny] {$t$};
				\draw[->] (0,0) -- (0,4)
				node[above,font=\tiny] {$U(t)$};
				% Mean line
				\draw[aerogreen,thick,dashed] (0,2) -- (5.2,2);
				\node[aerogreen,left,font=\tiny] at (0,2) {$\bar{U}$};
				% Signal
				\draw[gray,thick]
				(0.1,2.0) sin (0.4,2.6) cos (0.7,2.0) sin (1.0,1.4)
				cos (1.3,2.0) sin (1.6,2.8) cos (1.9,2.0)
				sin (2.2,1.3) cos (2.5,2.0) sin (2.8,2.5)
				cos (3.1,2.0) sin (3.4,1.5) cos (3.7,2.0)
				sin (4.0,2.7) cos (4.3,2.0) sin (4.6,1.6)
				cos (4.9,2.0);
				% RMS
				\draw[<->,aerored,thick] (5.0,2.0) -- (5.0,2.7);
				\node[aerored,right,font=\tiny] at (5.0,2.35) {$\sigma_U$};
			\end{tikzpicture}

			\vspace{0.3em}
			\centering\scriptsize
			\begin{tabular}{lr}
				\toprule
				\textbf{Carpeta} & \textbf{$I_u$ (\%)} \\
				\midrule
				FX01G00          & 11,68               \\
				FX02G00          & 8,94                \\
				FX03G00          & 8,08                \\
				FX04G00          & 8,35                \\
				\bottomrule
			\end{tabular}

			\begin{alertblock}{}
				\centering\scriptsize
				$I_u \approx 8$--$12\%$: turbulento moderado
			\end{alertblock}
		\end{column}
	\end{columns}
\end{frame}

% =============================================================================
\section{Resultados Experimentales}
% =============================================================================

\begin{frame}{Perfil de Velocidades $U(y)$}
	\centering
	\includegraphics[width=0.82\textwidth]{perfil_de_velocidades_anemometra_lser_doppler_lda.png}

	\vspace{0.3em}
	\scriptsize
	Perfil medido con LDA: $U$ crece con $y$ de
	\SI{0,52}{\m\per\s} (FX01) a \SI{2,98}{\m\per\s} (FX04)\\
	Barras de error = error estándar (SE $\sim\SI{0,001}{\m\per\s}$)
\end{frame}

\begin{frame}{Ajuste de Perfil Polinómico}
	\centering
	\includegraphics[width=0.82\textwidth]{perfil_ajuste_polinomico.png}

	\vspace{0.3em}
	\scriptsize
	Ajuste polinómico (grado~2):
	$U(y) = 2{,}97\!\times\!10^{-4}\,y^2 + 1{,}81\!\times\!10^{-2}\,y
		+ 0{,}283$ [m/s, mm]
\end{frame}

\begin{frame}{Intensidad de Turbulencia $I_u(y)$}
	\centering
	\includegraphics[width=0.82\textwidth]{perfil_intensidad_turbulencia_lda.png}

	\vspace{0.3em}
	\scriptsize
	$I_u$ mayor cerca de la pared (FX01: 11,7\%) y más uniforme
	en la zona central ($\sim$~8\%)
\end{frame}

\begin{frame}{Histogramas de Velocidad por Posición}
	\centering
	\includegraphics[width=0.82\textwidth]{histogramas_velocidad_U_limpia.png}

	\vspace{0.3em}
	\scriptsize
	Distribuciones aproximadamente gaussianas tras filtrado IQR.
	Nota la escala creciente en $U$ de FX01 a FX04.
\end{frame}

\begin{frame}{Box Plot -- Distribución de $U$ por Posición}
	\centering
	\includegraphics[width=0.75\textwidth]{boxplot_velocidades_U.png}

	\vspace{0.3em}
	\scriptsize
	Dispersión creciente con la velocidad media:
	$\sigma_U$ escala con $\bar{U}$ ($I_u$ aproximadamente constante)
\end{frame}

\begin{frame}{Coeficiente de Variación (CV) por Posición}
	\centering
	\includegraphics[width=0.75\textwidth]{CV_por_posicion.png}

	\vspace{0.3em}
	\scriptsize
	FX01 ligeramente por encima de CV~$=10\%$ (calidad regular);
	FX02--FX04 buena calidad (CV~$< 10\%$)
\end{frame}

\begin{frame}{Esfuerzo de Reynolds $\tau_\mathrm{turb}(y)$}
	\begin{columns}
		\begin{column}{0.55\textwidth}
			\centering
			\includegraphics[width=\textwidth]{perfil_esfuerzo_reynolds.png}
		\end{column}
		\begin{column}{0.42\textwidth}
			\textbf{Esfuerzo turbulento:}
			\begin{equation*}
				\tau_\mathrm{turb} = -\rho\,\overline{u'\,v'}
			\end{equation*}

			\begin{exampleblock}{Resultados}
				\scriptsize
				\begin{tabular}{lr}
					\toprule
					     & {$\tau$ (Pa)}   \\
					\midrule
					FX01 & $-\num{1,2e-4}$ \\
					FX02 & $-\num{5,7e-4}$ \\
					FX03 & $-\num{3,2e-4}$ \\
					FX04 & $-\num{6,9e-3}$ \\
					\bottomrule
				\end{tabular}
			\end{exampleblock}
			\vspace{0.2em}
			\scriptsize
			$R_{uv}$ bajo ($0{,}02$--$0{,}18$): correlación
			débil entre $u'$ y $v'$
		\end{column}
	\end{columns}
\end{frame}

\begin{frame}{Datos Brutos -- Archivo Individual}
	\centering
	\includegraphics[width=0.82\textwidth]{box_plot_de_velocidades_cols_4_y_5.png}

	\vspace{0.3em}
	\scriptsize
	Análisis de archivo \texttt{L01G00\_001.txt} (2000 muestras):
	$\bar{U} = \SI{0,377}{\m\per\s}$, $\bar{V} \approx 0$
\end{frame}

% =============================================================================
\section{Incertidumbre}
% =============================================================================

\begin{frame}{Propagación GUM -- Velocidad LDA}
	\textbf{Incertidumbre combinada:}
	\begin{equation*}
		u_U^2 = \delta_f^2\,u_{f_D}^2 + f_D^2\,u_{\delta_f}^2
	\end{equation*}

	\vspace{0.3em}
	\begin{columns}
		\begin{column}{0.5\textwidth}
			\textbf{Contribución dominante:} $u_{f_D}$

			\vspace{0.3em}
			BSA resolución $\Delta f = \SI{10}{\kHz}$\\
			(tipo~B, rectangular):
			\begin{equation*}
				u_{f_D} = \frac{\Delta f}{\sqrt{3}}
				\approx \SI{5,77}{\kHz}
			\end{equation*}

			\begin{equation*}
				u_U = \delta_f \cdot u_{f_D}
				\approx \SI{0,019}{\m\per\s}
			\end{equation*}
		\end{column}
		\begin{column}{0.47\textwidth}
			\textbf{Incertidumbre tipo~A:}
			\begin{equation*}
				u_{\bar{U}} = \frac{\sigma_U}{\sqrt{N}}
				\approx \SI{0,001}{\m\per\s}
			\end{equation*}
			($\sigma_U \approx 0{,}2$, $N \approx 30\,000$)

			\vspace{0.3em}
			\begin{block}{Incertidumbre expandida}
				\begin{align*}
					u_{c,U}        & \approx \SI{0,02}{\m\per\s} \\
					U_\mathrm{exp} & = 2\,u_{c,U}
					\approx \SI{0,04}{\m\per\s}                  \\
					u_U / U        & \approx 0{,}95\% \text{ a }
					\SI{2}{\m\per\s}
				\end{align*}
			\end{block}
		\end{column}
	\end{columns}
\end{frame}

% =============================================================================
\section{Caso de Estudio: Estela NACA~0012}
% =============================================================================

\begin{frame}{Coeficiente de Arrastre por Estela}
	\begin{block}{Método de Jones}
		\begin{equation*}
			C_D = \frac{2}{c}\int_{-\infty}^{+\infty}
			\frac{U_1}{U_\infty}\!\left(1
			- \frac{U_1}{U_\infty}\right)dy
			\approx 0{,}012
		\end{equation*}
		NACA~0012, $Re_c = 1{,}5\times10^5$, $\alpha = 0°$
		(consistente con literatura)
	\end{block}
	\begin{columns}
		\begin{column}{0.5\textwidth}
			\begin{itemize}\small
				\item $c = \SI{150}{\mm}$, medida a $x/c = 1{,}0$
				\item Barrido $y \in [-50,50]\;\si{\mm}$,
				      $\Delta y = \SI{10}{\mm}$
				\item Mínimo $U_1/U_\infty \approx 0{,}72$
				      en la línea central
				\item Integración trapezoidal con datos LDA
			\end{itemize}
		\end{column}
		\begin{column}{0.47\textwidth}
			\begin{tikzpicture}[scale=0.75,>=Stealth]
				\draw[->] (0,-2.5) -- (0,2.5)
				node[above,font=\tiny] {$y$};
				\draw[->] (-0.3,0) -- (4.0,0)
				node[right,font=\tiny] {$U/U_\infty$};
				% Wake profile
				\draw[aerogreen,very thick]
				plot[domain=-2.3:2.3,samples=40]
				({3.0 - 1.0*exp(-(\x)^2/0.7)}, \x);
				% Free stream
				\draw[gray,dashed] (3.0,-2.5) -- (3.0,2.5);
				% Deficit area
				\fill[aerogreen!15]
				plot[domain=-2.3:2.3,samples=40]
				({3.0-1.0*exp(-(\x)^2/0.7)},\x)
				-- (3.0,2.3) -- (3.0,-2.3) -- cycle;
				\node[aerogreen,font=\tiny] at (2.5,0) {Déficit};
				% Data points
				\foreach \y in {-2,-1.5,-1,-0.5,0,0.5,1,1.5,2} {
						\fill[aerored]
						({3.0-exp(-\y*\y/0.7)},\y) circle (0.07);
					}
				\node[gray,right,font=\tiny] at (3.0,2.3) {$U_\infty$};
			\end{tikzpicture}
		\end{column}
	\end{columns}
\end{frame}

% =============================================================================
\section{Evaluación de Riesgos}
% =============================================================================

\begin{frame}{Riesgos de la Práctica LDA}
	\centering\small
	\begin{tabular}{>{\raggedright\arraybackslash}p{2.6cm}
		c c c c l}
		\toprule
		\textbf{Riesgo} & \textbf{P}                             & \textbf{S}
		                & $\boldsymbol{R}$                       & \textbf{Nivel}
		                & \textbf{Controles}                                           \\
		\midrule
		\rowcolor{aerored!15}
		Exposición ocular al láser
		                & 3                                      & 5              & 15
		                & \textcolor{aerored}{\textbf{Alto}}
		                & Gafas OD${\geq}3$                                            \\
		\rowcolor{aerored!15}
		Reflexiones especulares
		                & 3                                      & 4              & 12
		                & \textcolor{aerored}{\textbf{Alto}}
		                & Carcasa cerrada                                              \\
		\rowcolor{aeroyellow!15}
		Fallo alineación óptica
		                & 3                                      & 3              & 9
		                & \textcolor{aeroyellow}{\textbf{Medio}}
		                & Recalibrar sesión                                            \\
		\rowcolor{aeroyellow!15}
		Inhalación seeding
		                & 2                                      & 3              & 6
		                & \textcolor{aeroyellow}{\textbf{Medio}}
		                & Ventilación                                                  \\
		\rowcolor{aerogreen!15}
		Descarga eléctrica
		                & 1                                      & 4              & 4
		                & \textcolor{aerogreen}{\textbf{Bajo}}
		                & Puesta a tierra                                              \\
		\rowcolor{aerogreen!15}
		Daño ventana óptica
		                & 2                                      & 2              & 4
		                & \textcolor{aerogreen}{\textbf{Bajo}}
		                & Paños ópticos                                                \\
		\bottomrule
	\end{tabular}

	\vspace{0.3em}
	\begin{block}{$R = \text{Probabilidad} \times \text{Severidad}$}
		Bajo ($<5$) \quad | \quad Medio (5--12)
		\quad | \quad Alto ($>12$)
	\end{block}
\end{frame}

\begin{frame}{Matriz de Riesgos $5\times5$}
	\centering
	\begin{tikzpicture}[scale=0.6]
		\foreach \x in {0,...,4} {
				\foreach \y in {0,...,4} {
						\pgfmathsetmacro{\risk}{(\x+1)*(\y+1)}
						\pgfmathparse{\risk < 5 ? "aerogreen!40" :
							(\risk < 13 ? "aeroyellow!60" : "aerored!50")}
						\edef\fillcol{\pgfmathresult}
						\fill[\fillcol] (\x,\y) rectangle (\x+1,\y+1);
						\draw[gray] (\x,\y) rectangle (\x+1,\y+1);
						\node[font=\tiny]
						at (\x+0.5,\y+0.5)
						{\pgfmathprintnumber[fixed,precision=0]{\risk}};
					}
			}
		\foreach \x/\l in
			{0/MuyBajo,1/Bajo,2/Medio,3/Alto,4/MuyAlto} {
				\node[font=\tiny,rotate=45,anchor=west]
				at (\x+0.5,-0.1) {\l};
			}
		\foreach \y/\l in
			{0/Insignif.,1/Menor,2/Moder.,3/Mayor,4/Catastrófico} {
				\node[font=\tiny,anchor=east] at (-0.1,\y+0.5) {\l};
			}
		\node[font=\tiny\bfseries,rotate=90] at (-1.3,2.5)
		{Severidad};
		\node[font=\tiny\bfseries] at (2.5,-1) {Probabilidad};
		% Risk markers
		\fill[blue,opacity=0.8] (2.5,4.5) circle (0.2);
		\node[white,font=\tiny\bfseries] at (2.5,4.5) {1};
		\fill[blue!70,opacity=0.8] (2.5,3.5) circle (0.2);
		\node[white,font=\tiny\bfseries] at (2.5,3.5) {2};
		\fill[orange,opacity=0.8] (2.5,2.5) circle (0.2);
		\node[white,font=\tiny\bfseries] at (2.5,2.5) {3};
		\fill[yellow!70,opacity=0.8] (1.5,2.5) circle (0.2);
		\node[font=\tiny\bfseries] at (1.5,2.5) {4};
		\fill[aerogreen!70,opacity=0.8] (0.5,3.5) circle (0.2);
		\node[white,font=\tiny\bfseries] at (0.5,3.5) {5};
		\fill[aerogreen!70,opacity=0.8] (1.5,1.5) circle (0.2);
		\node[white,font=\tiny\bfseries] at (1.5,1.5) {6};
	\end{tikzpicture}

	\vspace{0.3cm}
	\begin{tikzpicture}[scale=0.65,every node/.style={font=\small}]
		\fill[aerogreen!80] (0,0) -- (5,0) -- (4.2,0.8)
		-- (0.8,0.8) -- cycle;
		\fill[aerogreen!60] (0.8,0.8) -- (4.2,0.8) -- (3.4,1.6)
		-- (1.6,1.6) -- cycle;
		\fill[aeroyellow!80] (1.6,1.6) -- (3.4,1.6) -- (2.8,2.4)
		-- (2.2,2.4) -- cycle;
		\fill[aerored!60] (2.2,2.4) -- (2.8,2.4) -- (2.5,3.0) -- cycle;
		\node[white,font=\tiny\bfseries] at (2.5,0.4)
		{Eliminación / Sustitución};
		\node[white,font=\tiny\bfseries] at (2.5,1.2) {Ingeniería};
		\node[font=\tiny\bfseries] at (2.5,2.0) {Admin.};
		\node[white,font=\tiny\bfseries] at (2.5,2.6) {EPP};
		\draw[->,gray,very thick] (5.3,2.8) -- (5.3,-0.1);
		\node[gray,font=\tiny,rotate=90] at (5.6,1.3) {Más efectivo};
	\end{tikzpicture}
\end{frame}

% =============================================================================
\section{Casos de Estudio por Equipos}
% =============================================================================

\begin{frame}{Trabajo por Equipos: 4 Escenarios}
	\begin{block}{Instrucciones}
		Cada equipo recibe un escenario diferente.\ Deben:
		\begin{enumerate}
			\item Procesar los datos de la(s) carpeta(s) indicadas
			\item Calcular estadísticos: $\bar{U}$, $\sigma_U$,
			      $I_u$ por posición
			\item Construir el perfil $U(y)$ o comparar condiciones
			\item Evaluar riesgo: $P$ (1--5), $S$ (1--5),
			      $R = P\!\times\!S$, ubicar en la matriz
			\item Resolver los problemas adicionales del escenario
			\item Verificar resultados con el script Python
		\end{enumerate}
	\end{block}
\end{frame}

\begin{frame}{Equipo~1: Perfil de Velocidad Baja (FX01G00)}
	\begin{columns}
		\begin{column}{0.38\textwidth}
			\centering\scriptsize
			\begin{tabular}{ll}
				\toprule
				\textbf{Par.}  & \textbf{Valor}           \\
				\midrule
				Carpeta        & FX01G00                  \\
				$\bar{U}$ esp. & $\sim\SI{0,5}{\m\per\s}$ \\
				Archivos       & 15 $\times$ 2\,000       \\
				Criterio       & $I_u < 10\%$             \\
				\bottomrule
			\end{tabular}
		\end{column}
		\begin{column}{0.28\textwidth}
			\begin{alertblock}{\scriptsize Riesgo}
				\scriptsize
				Baja velocidad,\\
				señal débil,\\
				ruido del BSA.\\
				Evaluar $P$, $S$, $R$.
			\end{alertblock}
		\end{column}
		\begin{column}{0.3\textwidth}
			\begin{exampleblock}{\scriptsize Problemas}
				\scriptsize
				\begin{enumerate}
					\item[A.1] $\bar{U}$, $\sigma_U$, $I_u$
					      antes/después IQR
					\item[A.2] $f_D = \bar{U}/\delta_f$
					\item[A.3] $u_U$ y $u_U/\bar{U}$
				\end{enumerate}
			\end{exampleblock}
			\vspace{0.2em}
			\begin{tikzpicture}[scale=0.45,>=Stealth]
				\draw[fill=aerogreen!10,draw=aerogreen,thick]
				(0,0) rectangle (2.2,1.2);
				\draw[aerolaser,ultra thick] (0.5,0.6) -- (1.7,0.6);
				\foreach \y in {0.3,0.6,0.9} {
						\draw[->,blue!50,thick] (-0.3,\y) -- (0.2,\y);
					}
				\node[font=\tiny] at (1.1,1.0) {LDA};
			\end{tikzpicture}
		\end{column}
	\end{columns}
\end{frame}

\begin{frame}{Equipo~2: Comparación FX01 vs.\ FX04}
	\begin{columns}
		\begin{column}{0.38\textwidth}
			\centering\scriptsize
			\begin{tabular}{ll}
				\toprule
				\textbf{Par.}  & \textbf{Valor}           \\
				\midrule
				Carpetas       & FX01G00 y FX04G00        \\
				$\bar{U}$ FX01 & $\sim\SI{0,5}{\m\per\s}$ \\
				$\bar{U}$ FX04 & $\sim\SI{2,9}{\m\per\s}$ \\
				Análisis       & Perfiles superpuestos    \\
				\bottomrule
			\end{tabular}
		\end{column}
		\begin{column}{0.28\textwidth}
			\begin{alertblock}{\scriptsize Riesgo}
				\scriptsize
				Dos condiciones de\\
				velocidad extremas,\\
				rango del BSA.\\
				Evaluar $P$, $S$, $R$.
			\end{alertblock}
		\end{column}
		\begin{column}{0.3\textwidth}
			\begin{exampleblock}{\scriptsize Problemas}
				\scriptsize
				\begin{enumerate}
					\item[B.1] Histogramas $U$\\
					      ¿Gaussianos?
					\item[B.2] $U/\bar{U}$\\
					      ¿Colapsan?
					\item[B.3] $-\rho\overline{u'v'}$\\
					      ¿Escala con $\bar{U}^2$?
				\end{enumerate}
			\end{exampleblock}
		\end{column}
	\end{columns}
\end{frame}

\begin{frame}{Equipo~3: Calidad de Señal (FX02G00)}
	\begin{columns}
		\begin{column}{0.38\textwidth}
			\centering\scriptsize
			\begin{tabular}{ll}
				\toprule
				\textbf{Par.}  & \textbf{Valor}           \\
				\midrule
				Carpeta        & FX02G00                  \\
				Particularidad & 1\,138 muest./arch.      \\
				$\bar{U}$ esp. & $\sim\SI{1,0}{\m\per\s}$ \\
				Análisis       & Convergencia estat.      \\
				\bottomrule
			\end{tabular}
		\end{column}
		\begin{column}{0.28\textwidth}
			\begin{alertblock}{\scriptsize Riesgo}
				\scriptsize
				Menor $N$ de muestras,\\
				posible bajo seeding,\\
				sesgo estadístico.\\
				Evaluar $P$, $S$, $R$.
			\end{alertblock}
		\end{column}
		\begin{column}{0.3\textwidth}
			\begin{exampleblock}{\scriptsize Problemas}
				\scriptsize
				\begin{enumerate}
					\item[C.1] ¿Por qué menos\\
					      muestras?
					\item[C.2] $\bar{U}$ acumulado\\
					      ¿Converge a $\pm1\%$?
					\item[C.3] $\sigma_U$ con/sin\\
					      filtrado IQR
				\end{enumerate}
			\end{exampleblock}
		\end{column}
	\end{columns}
\end{frame}

\begin{frame}{Equipo~4: Estela NACA~0012 (FX03 + FX04)}
	\begin{columns}
		\begin{column}{0.38\textwidth}
			\centering\scriptsize
			\begin{tabular}{ll}
				\toprule
				\textbf{Par.} & \textbf{Valor}         \\
				\midrule
				Carpetas      & FX03G00 y FX04G00      \\
				Perfil        & NACA~0012, $\alpha=0°$ \\
				$Re_c$        & $1{,}5\times10^5$      \\
				Referencia    & $C_D \approx 0{,}012$  \\
				\bottomrule
			\end{tabular}
		\end{column}
		\begin{column}{0.28\textwidth}
			\begin{alertblock}{\scriptsize Riesgo}
				\scriptsize
				Sonda cerca del perfil,\\
				reflexiones del borde\\
				de salida, seeding\\
				no uniforme. $P$, $S$, $R$.
			\end{alertblock}
		\end{column}
		\begin{column}{0.3\textwidth}
			\begin{exampleblock}{\scriptsize Problemas}
				\scriptsize
				\begin{enumerate}
					\item[D.1] Perfil $U/U_\infty$
					\item[D.2] $C_D$ por integración
					      trapezoidal
					\item[D.3] $\theta$ y comparar
					      con $C_D \cdot c / 2$
				\end{enumerate}
			\end{exampleblock}
			\vspace{0.2em}
			\begin{tikzpicture}[scale=0.45,>=Stealth]
				% NACA schematic
				\draw[thick,fill=gray!30]
				(0,0.08) .. controls (0.5,0.15) and (1.5,0.12)
				.. (2.2,0.02)
				-- (2.2,-0.02)
				.. controls (1.5,-0.12) and (0.5,-0.15)
				.. (0,-0.08) -- cycle;
				% Wake
				\draw[aerogreen,thick,dashed] (2.2,0.5) -- (3.0,0.5);
				\draw[aerogreen,thick,dashed] (2.2,-0.5) -- (3.0,-0.5);
				\node[font=\tiny] at (2.8,0) {Estela};
			\end{tikzpicture}
		\end{column}
	\end{columns}
\end{frame}

% =============================================================================
\section{Herramientas de Software}
% =============================================================================

\begin{frame}[fragile]{Script \texttt{P04\_LDA\_Perfil\_Velocidad.py}}
	\begin{columns}
		\begin{column}{0.5\textwidth}
			\begin{block}{Ruta}
				\scriptsize
				\texttt{Practicas/P04\_LDA\_Perfil\_Velocidad/}\\
				\texttt{\ \ src/P04\_LDA\_Perfil\_Velocidad.py}
			\end{block}
			\vspace{0.2em}
			\scriptsize
			\begin{tabular}{ll}
				\toprule
				\textbf{Función}
				 & \textbf{Resultado}                      \\
				\midrule
				\texttt{remove\_outliers\_iqr}
				 & Array filtrado + máscara                \\
				\texttt{process\_position\_folder}
				 & Dict: $\bar{U}$, $\sigma_U$, $I_u$, $N$ \\
				Bloque principal
				 & Perfiles, histogramas,                  \\
				 & box-plots, Reynolds                     \\
				\bottomrule
			\end{tabular}
			\vspace{0.3em}

			\textbf{Salidas:}
			\begin{itemize}\scriptsize
				\item \texttt{data/images/*.png} y \texttt{*.svg}
				\item 8 gráficas generadas
				\item Tabla resumen en consola
			\end{itemize}
		\end{column}
		\begin{column}{0.47\textwidth}
			\begin{lstlisting}[language=Python]
import numpy as np
import pandas as pd
from pathlib import Path
from scipy import stats

def remove_outliers_iqr(data, k=1.5):
    Q1 = np.percentile(data, 25)
    Q3 = np.percentile(data, 75)
    IQR = Q3 - Q1
    mask = (data >= Q1 - k*IQR)
         & (data <= Q3 + k*IQR)
    return data[mask], mask

def process_position_folder(folder):
    all_U = []
    for f in folder.glob('*.txt'):
        df = pd.read_csv(f, sep=r'\s+',
            header=None,
            names=['idx','t','SNR','U','V'])
        all_U.extend(df['U'].values)
    U = np.array(all_U)
    U_clean, _ = remove_outliers_iqr(U)
    return {
        'U_mean': np.mean(U_clean),
        'U_std': np.std(U_clean),
        'I_u': np.std(U_clean)
             / np.mean(U_clean) * 100,
        'N': len(U_clean)
    }
			\end{lstlisting}
		\end{column}
	\end{columns}
\end{frame}

\begin{frame}{Flujo de Ejecución del Script}
	\begin{columns}
		\begin{column}{0.5\textwidth}
			\textbf{\texttt{P04\_LDA\_Perfil\_Velocidad.py}}\\[0.3em]
			\scriptsize
			\begin{enumerate}
				\item Descubrimiento de \texttt{files/P4/}
				      (4 carpetas)
				\item Lectura de 60 archivos \texttt{.txt}
				\item Concatenación $U$ y $V$ por posición
				\item Filtrado IQR ($k = 1{,}5$)
				\item Estadísticos: $\bar{U}$, $\sigma_U$,
				      $I_u$, SE
				\item Perfil $U(y)$ con barras de error
				\item Histogramas por posición
				\item Box-plots y coef.\ de variación
				\item Ajuste polinómico (grado~2)
				\item Esfuerzos de Reynolds
				\item Gráficas PNG + SVG
			\end{enumerate}
		\end{column}
		\begin{column}{0.47\textwidth}
			\begin{exampleblock}{Ejecución}
				\texttt{cd Practicas/.../src}\\
				\texttt{python P04\_LDA\_...py}
			\end{exampleblock}
			\vspace{0.3em}
			\begin{block}{Gráficas generadas}
				\scriptsize
				\begin{itemize}
					\item \texttt{perfil\_de\_velocidades\_*.png}
					\item \texttt{perfil\_intensidad\_*.png}
					\item \texttt{histogramas\_velocidad\_*.png}
					\item \texttt{boxplot\_velocidades\_*.png}
					\item \texttt{CV\_por\_posicion.png}
					\item \texttt{perfil\_ajuste\_*.png}
					\item \texttt{perfil\_esfuerzo\_*.png}
					\item \texttt{box\_plot\_de\_velocidades\_*.png}
				\end{itemize}
			\end{block}
		\end{column}
	\end{columns}
\end{frame}

% =============================================================================
\section{Cuestionario y Formato}
% =============================================================================

\begin{frame}{Preguntas de Discusión}
	\begin{enumerate}\small
		\item ¿Cuál es la magnitud medida directamente por el sistema
		      LDA? ¿Cómo se diferencia de la HWA?
		\item Deducir $\delta_f$ a partir de los parámetros ópticos.
		      ¿Qué ocurre si se reduce $\theta/2$?
		\item ¿Por qué la LDA no necesita calibración?
		\item Función de la celda de Bragg: ¿qué pasaría sin ella
		      en un flujo con recirculación?
		\item ¿Cómo afecta la densidad de seeding a la tasa de
		      adquisición?
		\item ¿Por qué el IQR ($k = 1{,}5$) es más robusto que
		      $\pm2\sigma$?
		\item Interpretar los valores de $I_u$: ¿laminar, transición
		      o turbulento?
		\item Si se usara un Pitot en vez de LDA, ¿cómo afectaría
		      la resolución cerca de la pared?
		\item Dos aplicaciones aeronáuticas donde LDA sea superior
		      al hilo caliente y justificar por qué
		\item ¿Qué información adicional proporciona
		      $-\rho\overline{u'v'}$?
		\item Criterio de Stokes ($\mathrm{Stk} \ll 1$) para
		      trazadores
	\end{enumerate}
\end{frame}

\begin{frame}{Formato del Informe (AIAA)}
	\begin{enumerate}
		\item \textbf{Resumen:} máx.\ 200 palabras
		\item \textbf{Introducción}
		\item \textbf{Fundamento Teórico:} efecto Doppler, modelo
		      de franjas, Bragg
		\item \textbf{Metodología:} procesamiento + herramientas
		\item \textbf{Resultados:} tabla de estadísticos por carpeta,
		      perfil $U(y)$, histogramas, caso del equipo,
		      evaluación de riesgo
		\item \textbf{Discusión:} cuestionario + análisis de
		      incertidumbre
		\item \textbf{Conclusiones:} 4--6 puntos
		\item \textbf{Referencias} (IEEE)
		\item \textbf{Apéndice~A:} gráficas Python
	\end{enumerate}
\end{frame}

% =============================================================================
\section{Conclusiones}
% =============================================================================

\begin{frame}{Conclusiones}
	\begin{enumerate}
		\item El sistema LDA Ar-Ion ($\lambda = \SI{514,5}{\nm}$) con
		      $\delta_f = \SI{3,28}{\micro\m}$ mide velocidades con
		      incertidumbre relativa $< 1\%$
		\item El volumen de medida
		      ($d_{mv} \approx \SI{137}{\micro\m}$,
		      $l_{mv} \approx \SI{1,74}{\mm}$) proporciona resolución
		      espacial adecuada para perfiles de capa límite
		\item El filtrado IQR ($k = 1{,}5$) elimina 3--6\% de
		      outliers por posición, mejorando la calidad estadística
		\item Las 4 condiciones muestran $\bar{U}$ de
		      \SIrange{0,5}{3,0}{\m\per\s} con $I_u \approx 8$--$12\%$
		      (flujo moderadamente turbulento)
		\item La componente transversal $\bar{V} < 3\%\,\bar{U}$
		      confirma flujo predominantemente unidimensional
		\item La integración del déficit de estela permite estimar
		      $C_D \approx 0{,}012$ para NACA~0012
		\item El script Python automatiza lectura, filtrado,
		      estadísticos, perfiles y gráficas
	\end{enumerate}

	\vspace{0.3em}
	\centering
	\begin{tikzpicture}
		\node[draw=aerogreen,thick,rounded corners,fill=aerogreen!10,
			text width=10cm,align=center,font=\small] {
			\textit{<<\,La LDA es la técnica de referencia para medición
				no intrusiva de velocidad en flujos
				aerodinámicos\,>>}};
	\end{tikzpicture}
\end{frame}

\begin{frame}{}
	\centering
	\vspace{2cm}
	{\Huge\color{aerogreen}\bfseries ¿Preguntas?}\\[1cm]
	{\large P04 -- Anemometría Láser Doppler (LDA)}\\[0.5cm]
	{\small Técnicas de Medida -- Ingeniería Aeronáutica}
\end{frame}

\end{document}
